% Section 3.3 -- retention and its scope decomposition.
% Characterization, not apology. No literal numbers: every figure is a macro
% from numbers.tex. Budget: 500-600 words (section 3 total is 2.25 pages).

\subsection{Retention, and what its denominator is a property of}
\label{sec:retention}

We report retention against a decomposition rather than against a threshold,
and the reason is instructive about how instruments in this setting should be
scored at all.

Our original design gated the project on retention: the extractor had to clear
a fixed fraction of the explicit targets in its sample or the approach was
abandoned. That gate was retired before it ever returned a verdict we acted on,
because its threshold turned out to encode two things that have nothing to do
with our parser. The first is the elaborator behaviour of
Section~\ref{sec:is-explicit}: part of the explicit set was never in the source
text, so no fraction of it is reachable by reading source. The second is
editorial: mathlib authors group related results in files, so whether a premise
name happens to appear elsewhere in a declaration's file is a fact about how the
library is organised. A threshold that moves when either of those changes is not
measuring the instrument, and a pass it awards is not evidence about the
instrument either. What replaces it is a decomposition into quantities that
are each measured the same way, so that the steps between them can be read
off rather than argued for.

\begin{table}[h]
\centering
\caption{Retention of mathematical-layer explicit targets, same sample and same
loose matcher, at three scopes. These are matcher-based reference points, not
attributable bounds: a loose name match is evidence that a name occurs within
the scope, not that the premise was used there, and it can equally miss a
premise written under another spelling or notation.}
\label{tab:retention}
\small
\begin{tabular}{llc}
\toprule
Quantity & Denominator & Value \\
\midrule
Whole-file reference   & explicit targets, name matched anywhere in the file & \ceilingWholeFile{} \\
Declaration reference  & explicit targets, name matched in the declaration   & \reachableBound{} \\
Achieved retention     & explicit targets, our extractor                     & \retention{} \\
Efficiency             & achieved against the declaration reference          & \efficiencyReachable{} \\
\bottomrule
\end{tabular}
\end{table}

Read downward, Table~\ref{tab:retention} separates three steps of very
different kinds. Targets missed even by a whole-file match are not recoverable
by reading source under any spelling the matcher knows; elaborator-inserted
constants are the mechanism we can name for that, and we do not claim it
accounts for all of them. The step from the whole-file to the declaration
reference is an organisational effect: those targets are named somewhere in the
file but nowhere in the declaration that depends on them, so a per-declaration
instrument cannot see them however good its name resolution is. Only the last
line is a claim about our parser, and it is the figure we ask to be judged on.

The \residualTargets{} targets the declaration reference point finds and the
extractor misses are identified rather than merely counted, and
Appendix~\ref{app:residual} decomposes them: the dominant class is the
definitional gap between resolving an identifier and matching a substring, and
closing it entirely would mean buying retention at the cost of precision. Two
properties make this a characterization rather than a disclosure. Each residual
class names a mechanism. And efficiency is close to flat across tactic classes,
from \efficiencyByClassMin{} to \efficiencyByClassMax{}, so the instrument does
not fail preferentially on the proof styles that distinguish the two corpora.
That is a statement about mathlib proofs only, as
Section~\ref{sec:limitations}, L1, records.
